\section{Real Evidence Package}

The active package is
\texttt{docs/examples/g-ensemble-evidence-packages/G.1-G.3+real-2020/}.
It binds the BISECT baselines, 2016 and 2020 election inputs, Rust and
GerryChain traces, RPLAN/RCTX contexts, audit certificates, analysis, and
verification commands.

\subsection{Design}

The preregistered design used four independent chains and 2,000 steps per chain
at 0.5\% population tolerance, discarding the first 500 steps. Partition
snapshots were retained every ten steps. Rust used a random-weight Kruskal tree
kernel matching GerryChain 0.3.2 for this comparison
~\citep{deford2021recombination,gerrychain2021}. The default Wilson UST
kernel remains a distinct sampler and is not conflated with this cross-tool
experiment.

Wisconsin failed the input eligibility gate: district 1 was disconnected under
the sampled adjacency graph. The failure remains archived. Iowa was added as a
medium-state replacement before any production chain metric was inspected.

\subsection{Cut-Fraction Results}

\begin{table}[h]
\centering
\caption{Benchmark percentile in the unweighted tract-graph cut-fraction distribution.}
\begin{tabular}{lrrrrrr}
\toprule
State & Tool & Percentile & 95\% CI & $\widehat R$ & ESS & Accept. \\
\midrule
RI & Rust & 1.35 & [0.93, 1.77] & 1.0005 & 4290 & 0.790 \\
RI & GerryChain & 1.81 & [1.48, 2.07] & 1.0000 & 6000 & 1.000 \\
IA & Rust & 0.33 & [0.20, 0.47] & 1.0046 & 257 & 0.997 \\
IA & GerryChain & 0.88 & [0.46, 1.18] & 1.0005 & 498 & 0.997 \\
NC & Rust & 0.00 & [0.00, 0.00] & 1.0404 & 109 & 1.000 \\
NC & GerryChain & 0.20 & [0.00, 0.40] & 1.0070 & 87 & 0.993 \\
\bottomrule
\end{tabular}
\end{table}

Rhode Island supports a bounded low-tail result in both tools. Iowa suggests a
low-tail benchmark, but neither ESS reaches the preregistered 1,000 threshold
for a sub-1\% headline. North Carolina cannot support a precise tail claim:
GerryChain ESS is below 100 and Rust is only slightly above it.

\subsection{Cross-Tool Differences}

The cut distributions are not interchangeable. Two-sample Kolmogorov--Smirnov
statistics are 0.038 for Rhode Island, 0.162 for Iowa, and 0.471 for North
Carolina. With thousands of autocorrelated draws, the associated nominal
$p$-values are not independent-sample legal tests; the effect sizes nevertheless
show material kernel/implementation sensitivity in Iowa and especially North
Carolina.

Partition-space integrated autocorrelation times from ten-step snapshots are
approximately 21 for Rhode Island, 11 for Iowa, and 7--8 for North Carolina.
Scalar convergence alone therefore does not imply independent partitions.

\section{Claim Boundary}

Cut fraction is a topology proxy, not Polsby--Popper or Reock. These three
tract-level states do not establish a national compactness extremum, neutral
distribution, legal validity, or block-level NRS conformance. The former
Wisconsin, Georgia, Pennsylvania, Texas, California, and North Carolina
percentile claims are withdrawn unless regenerated by an active package.

\section{Conclusion}

The strongest result is methodological. A benchmark percentile is not one
number: it depends on graph identity, tolerance, proposal kernel, chain
diagnostics, and implementation. Rhode Island supplies a reproducible
low-tail example; Iowa and North Carolina supply cautionary evidence against
unqualified extremum claims.
